% ==============================================================================
% Phân mục 00: Tóm tắt bài báo có cấu trúc (Structured Abstract chuẩn JMIR/TRIPOD-AI)
% ==============================================================================
\begin{clinicalbox}{TÓM TẮT NGHIÊN CỨU (STRUCTURED ABSTRACT --- CHUẨN TRIPOD-AI)}
\textbf{Bối cảnh (Background):} Các mô hình trí tuệ nhân tạo hỗ trợ chẩn đoán lâm sàng từ Hồ sơ sức khỏe điện tử (EHR) hiện nay chủ yếu dựa vào dữ liệu bảng có cấu trúc và bỏ qua nguồn ngữ cảnh phong phú trong ghi chú bệnh án tự do. Theo tổng quan hệ thống trên tạp chí JMIR (2024), 100\% (15/15) nghiên cứu đạt chuẩn sàng lọc chỉ sử dụng dữ liệu bảng và không có nghiên cứu nào đạt mức nguy cơ sai lệch thấp theo tiêu chuẩn PROBAST. \\[0.4em]
\textbf{Mục tiêu (Objective):} Nghiên cứu, thiết kế và hiện thực hóa nguyên mẫu Hệ thống Hỗ trợ Quyết định Lâm sàng Đa phương thức (M-CDSS) kết hợp dữ liệu sinh hiệu/xét nghiệm dạng bảng với văn bản bệnh án tự do thông qua cơ chế chú ý chéo (\textit{Cross-Attention Late Fusion}), tích hợp hệ thống đa tác tử truy xuất y văn (\textit{Agentic RAG}) có kiểm chứng vị trí ký tự chống ảo giác và hiệu chuẩn xác suất đạt chuẩn lâm sàng. \\[0.4em]
\textbf{Phương pháp (Methods):} Nghiên cứu được thực hiện trên bộ dữ liệu 2.000 hồ sơ bệnh án đa phương thức cân bằng giữa 4 bệnh lý hô hấp -- tim mạch cấp tính (\texttt{ASTHMA}, \texttt{COPD}, \texttt{HEART\_FAILURE}, \texttt{PNEUMONIA}), phân chia nghiêm ngặt theo mã định danh bệnh nhân (\textit{Strict Group Split}, tập kiểm thử độc lập $N=302$, 0\% rò rỉ dữ liệu). Nhánh bảng số sử dụng \textit{FT-Transformer}, nhánh văn bản sử dụng \textit{Bio\_ClinicalBERT}, dung hợp qua tầng \textit{Multi-Head Cross-Attention} và hiệu chuẩn hậu kỳ bằng \textit{Temperature Scaling}. Kiến trúc phần mềm được đóng gói dạng vi dịch vụ 3 tầng (\texttt{Next.js 14}, \texttt{FastAPI}, \texttt{Qdrant Vector DB}) vận hành qua giao thức RESTful và đồ thị trạng thái \texttt{DiagnosticStateGraph}. \\[0.4em]
\textbf{Kết quả (Results):} Trên tập kiểm thử độc lập phân tách nghiêm ngặt theo bệnh nhân ($N=302$), thực nghiệm Giai đoạn 1 trên 3 mô hình đơn phương thức ghi nhận \textit{CatBoost} đạt $\text{AUROC} = 0.9223$, \textit{FT-Transformer} đạt $\text{AUROC} = 0.8976$ ($\text{ECE} = 0.1152$) và \textit{Bio\_ClinicalBERT} đạt $\text{AUROC} = 0.9735$ ($\text{Macro-F1} = 0.8421$). Đối với mô hình dung hợp \textbf{M-CDSS} Giai đoạn 2: \texttt{[TODO-TV1/TV2 (Backlog 04 Task 1.4, 2.3 \& Backlog 06 Task 1.1): Điền các chỉ số thực nghiệm AUROC macro, AUPRC macro, Macro-F1, Brier Score, ECE sau hiệu chuẩn Temperature Scaling từ checkpoint multimodal\_best.pt]}. Phân tích đường cong quyết định (DCA): \texttt{[TODO-TV1/TV2 (Backlog 06 Task 1.2 \& 2.2): Điền giá trị Net Benefit thực nghiệm, dải ngưỡng hiệu quả lâm sàng $p_t$, số ca can thiệp đúng ròng $\Delta\text{TP}_{\text{net}/100}$ và số ca can thiệp thừa tránh được $\Delta\text{FP}_{\text{avoided}/100}$ sau khi xuất ma trận dự đoán hold-out]}. Bộ kiểm chứng \texttt{EvidenceGrounder} đạt tỷ lệ xác thực 100\% trích dẫn khớp nguyên văn bệnh án gốc trên bộ kiểm thử tự động. \\[0.4em]
\textbf{Kết luận (Conclusions):} M-CDSS chứng minh tính bổ trợ khoa học chặt chẽ giữa ``mỏ neo huyết động học'' của dữ liệu bảng và ngữ cảnh bệnh sử của dữ liệu văn bản, đồng thời thiết lập một kiến trúc phần mềm y tế minh bạch, tái lập được qua đóng gói container.
\end{clinicalbox}
